\clearpage
\phantomsection
\addcontentsline{toc}{chapter}{Simulado II: Revisão Integrada da Unidade II}
\chapter*{Simulado II: Revisão Integrada da Unidade II (Processo e Estratégia)}
\markboth{SIMULADO II, UNIDADE II}{SIMULADO II, UNIDADE II}

\noindent\tagAV{Base AV3} \tagQuant\ \tagTerritorio

\begin{tcolorbox}[colback=bgSoft, colframe=etIndigo, arc=3pt, boxrule=0.7pt, left=10pt, right=10pt, top=7pt, bottom=7pt]
  \textbf{Instruções para Estudo Autônomo:} Este simulado integra os conteúdos dos \textbf{Capítulos 7 a 11} (Unidade II, Processo e Estratégia). Resolva as \textbf{8 questões objetivas (A–E)} e as \textbf{2 questões dissertativas} sem consultar o gabarito. Ao concluir, confira o gabarito comentado, o espelho dissertativo e preencha a Tabela de Diagnóstico na página seguinte.
\end{tcolorbox}

\section*{Parte A: Questões Objetivas (Capítulos 7 a 11)}

\begin{enumerate}[leftmargin=1.5em, itemsep=8pt]
  \item \textbf{(Cap. 7, Processo Empreendedor e \textit{Effectuation})} Um grupo de pesquisadores da UEFS desenvolveu um sensor óptico para detectar adulteração em méis regionais. Sem histórico de vendas desse dispositivo, decidiram limitar o investimento inicial a R\$\,15.000,00 próprios e firmar pré-compromisso com duas cooperativas de apicultores que cederam amostras e espaço de teste em troca de prioridade no uso. À luz de Sarasvathy (2001), essa estratégia exemplifica, respectivamente, os princípios:
  \begin{alternativas}
    \item Causation preditiva e maximização do retorno esperado.
    \item Perda aceitável (\textit{affordable loss}) e colcha de retalhos (\textit{crazy quilt}).
    \item Limonada (\textit{lemonade}) e análise competitiva de Porter.
    \item Piloto no avião e planejamento causal de longo prazo.
    \item Nexo indivíduo–oportunidade e diluição acionária.
  \end{alternativas}
  \espacoResolucao[1.2cm]

  \item \textbf{(Cap. 7, Curva J e Vale da Morte Tecnológico)} Uma startup de biotecnologia apresenta os seguintes fluxos de caixa líquidos anuais (em R\$): Ano 0 = $-90.000$; Ano 1 = $-110.000$; Ano 2 = $-40.000$; Ano 3 = $+120.000$; Ano 4 = $+160.000$. O fundo do Vale da Morte e o prazo de \textit{payback} simples desse empreendimento são, respectivamente:
  \begin{alternativas}
    \item $-\text{R\$\,200.000,00}$ no Ano 1 e 3{,}50 anos.
    \item $-\text{R\$\,240.000,00}$ no Ano 2 e 3{,}75 anos.
    \item $-\text{R\$\,240.000,00}$ no Ano 2 e 3{,}25 anos.
    \item $-\text{R\$\,120.000,00}$ no Ano 3 e 4{,}00 anos.
    \item $-\text{R\$\,240.000,00}$ no Ano 1 e 2{,}75 anos.
  \end{alternativas}
  \espacoResolucao[1.5cm]

  \item \textbf{(Cap. 8, Dimensionamento TAM, SAM e SOM)} Uma empresa desenvolveu uma plataforma de gestão hídrica para pequenas agroindústrias. No Brasil existem $50.000$ agroindústrias desse porte, com gasto médio potencial de R\$\,2.400,00/ano. O recorte regional do Nordeste e perfil aderente representa 30\% desse universo. Sabendo que a equipe comercial tem capacidade operacional (\textit{bottom-up}) para atender no máximo $450$ agroindústrias por ano no horizonte de 3 anos, o TAM, o SAM e o SOM \textit{bottom-up} anuais são, respectivamente:
  \begin{alternativas}
    \item R\$\,120,0 milhões; R\$\,36,0 milhões; R\$\,1,08 milhão.
    \item R\$\,120,0 milhões; R\$\,18,0 milhões; R\$\,1,08 milhão.
    \item R\$\,36,0 milhões; R\$\,10,8 milhões; R\$\,1,08 milhão.
    \item R\$\,120,0 milhões; R\$\,36,0 milhões; R\$\,3,60 milhões.
    \item R\$\,60,0 milhões; R\$\,36,0 milhões; R\$\,0,90 milhão.
  \end{alternativas}
  \espacoResolucao[1.5cm]

  \item \textbf{(Cap. 8, Propriedade Intelectual e Anterioridade)} Uma startup baiana desenvolveu um método inédito de extração enzimática de óleo de licuri (processo industrial) e um algoritmo de controle de temperatura embarcado no extrator. Uma busca no INPI identificou que o processo enzimático não possui anterioridade mundial, mas o algoritmo de controle baseia-se em código aberto sob licença permissiva. A estratégia de proteção mais adequada para os ativos próprios da startup é:
  \begin{alternativas}
    \item Modelo de Utilidade para o algoritmo e Marca Coletiva para o processo enzimático.
    \item Patente de Invenção para o processo enzimático inédito e Segredo Industrial ou Registro de Software apenas para as adaptações proprietárias do controle, respeitando a licença do código-base.
    \item Indicação Geográfica para o processo enzimático e Cultivar para o extrator mecânico.
    \item Patente de Invenção para o código aberto de terceiros e Direito Autoral para o óleo de licuri.
    \item Nenhuma proteção é possível, pois o uso de qualquer componente de código aberto impede o depósito de patente do processo químico.
  \end{alternativas}
  \espacoResolucao[1.2cm]

  \item \textbf{(Cap. 9, \textit{Business Model Canvas} e Economia Unitária)} Um serviço de monitoramento de solo por assinatura cobra $\text{ARPU} = \text{R\$\,400,00/mês}$, opera com margem bruta de 75\% e apresenta taxa mensal de \textit{churn} de 5\%. Se a empresa gastou R\$\,15.000,00 em marketing e vendas para adquirir 10 novos clientes no mês, o CAC, o LTV e a razão LTV/CAC são, respectivamente:
  \begin{alternativas}
    \item $\text{CAC} = \text{R\$\,1.500,00}$; $\text{LTV} = \text{R\$\,8.000,00}$; $\text{LTV/CAC} = 5{,}33$.
    \item $\text{CAC} = \text{R\$\,1.500,00}$; $\text{LTV} = \text{R\$\,6.000,00}$; $\text{LTV/CAC} = 4{,}00$.
    \item $\text{CAC} = \text{R\$\,1.000,00}$; $\text{LTV} = \text{R\$\,6.000,00}$; $\text{LTV/CAC} = 6{,}00$.
    \item $\text{CAC} = \text{R\$\,1.500,00}$; $\text{LTV} = \text{R\$\,4.500,00}$; $\text{LTV/CAC} = 3{,}00$.
    \item $\text{CAC} = \text{R\$\,3.000,00}$; $\text{LTV} = \text{R\$\,6.000,00}$; $\text{LTV/CAC} = 2{,}00$.
  \end{alternativas}
  \espacoResolucao[1.5cm]

  \item \textbf{(Cap. 10, Escada do Capital e Indicadores de Tração)} Uma empresa de base tecnológica possui caixa disponível de R\$\,320.000,00, despesas operacionais mensais totais (\textit{Gross Burn}) de R\$\,60.000,00 e 50 clientes pagando assinatura mensal de R\$\,400,00. Considerando essas condições constantes, o MRR, o \textit{Net Burn} e o \textit{Runway} estático são, respectivamente:
  \begin{alternativas}
    \item $\text{MRR} = \text{R\$\,20.000,00}$; $\text{Net Burn} = \text{R\$\,40.000,00/mês}$; $\text{Runway} = 8{,}0\text{ meses}$.
    \item $\text{MRR} = \text{R\$\,20.000,00}$; $\text{Net Burn} = \text{R\$\,60.000,00/mês}$; $\text{Runway} = 5{,}3\text{ meses}$.
    \item $\text{MRR} = \text{R\$\,25.000,00}$; $\text{Net Burn} = \text{R\$\,35.000,00/mês}$; $\text{Runway} = 9{,}1\text{ meses}$.
    \item $\text{MRR} = \text{R\$\,20.000,00}$; $\text{Net Burn} = \text{R\$\,40.000,00/mês}$; $\text{Runway} = 16{,}0\text{ meses}$.
    \item $\text{MRR} = \text{R\$\,30.000,00}$; $\text{Net Burn} = \text{R\$\,30.000,00/mês}$; $\text{Runway} = 10{,}7\text{ meses}$.
  \end{alternativas}
  \espacoResolucao[1.5cm]

  \item \textbf{(Cap. 10, Fomento Regional e Diluição)} Assinale a alternativa que apresenta o alinhamento correto entre estágio tecnológico e instrumento da Escada do Capital para uma EBT localizada na Bahia:
  \begin{alternativas}
    \item TRL 2–3 deve ser financiado prioritariamente por crédito bancário comercial de curto prazo com garantia real dos sócios.
    \item Subvenções econômicas como Programa Centelha, FAPESB e FINEP Tecnova têm caráter não dilutivo e compartilham o risco tecnológico na travessia TRL 3–6.
    \item O FNE Inovação do Banco do Nordeste atua exclusivamente como capital de risco (\textit{equity}), exigindo assento no conselho e cessão de 20\% das quotas.
    \item Investidores-anjo substituem a necessidade de validação de MVP, pois aportam capital apenas em ideias sem contato prévio com clientes.
    \item O \textit{Pitch Deck} substitui integralmente o Plano de Negócio Completo na aprovação de linhas de financiamento industrial do BNDES.
  \end{alternativas}
  \espacoResolucao[1.2cm]

  \item \textbf{(Cap. 11, \textit{Open Innovation}, \textit{Spin-offs} e Lei nº 13.243/2016)} Uma \textit{spin-off} acadêmica licenciou junto ao NIT de uma universidade pública uma patente de biofilme para frutas sob alíquota de \textit{royalty} de 5\% sobre o faturamento líquido. No ano de referência, a empresa obteve faturamento líquido de R\$\,1.000.000,00 com o produto. Sabendo que a resolução interna da ICT destina 30\% dos ganhos econômicos aos inventores (dentro do intervalo legal de 5\% a 1/3 da Lei nº 10.973/2004), o valor anual transferido aos inventores e o valor retido pela ICT/NIT são, respectivamente:
  \begin{alternativas}
    \item R\$\,50.000,00 e R\$\,100.000,00.
    \item R\$\,15.000,00 e R\$\,35.000,00.
    \item R\$\,16.666,67 e R\$\,33.333,33.
    \item R\$\,25.000,00 e R\$\,25.000,00.
    \item R\$\,5.000,00 e R\$\,45.000,00.
  \end{alternativas}
  \espacoResolucao[1.5cm]
\end{enumerate}

\section*{Parte B: Questões Dissertativas (Integração Estratégica e Territorial)}

\begin{enumerate}[leftmargin=1.5em, itemsep=8pt, start=9]
  \item \textbf{(Integração Caps. 7, 8 e 9, Do Vale da Morte ao Modelo de Negócio)} Um laboratório universitário desenvolveu um secador solar automatizado para amêndoas de cacau (TRL 4), mas não encontra compradores dispostos a pagar R\$\,35.000,00 à vista pelo equipamento. Proponha uma estratégia integrada de travessia do Vale da Morte combinando: (a) dois princípios de \textit{Effectuation}; (b) cruzamento de SOM \textit{bottom-up}; e (c) reformulação do bloco Fontes de Receita do \textit{Business Model Canvas} (Chesbrough \& Rosenbloom, 2002).
  \espacoResolucao[2.8cm]

  \item \textbf{(Integração Caps. 10 e 11, Escada do Capital e \textit{Spin-off} Acadêmica)} Explique como uma \textit{spin-off} acadêmica originada em uma ICT pública pode combinar a Lei nº 13.243/2016 (compartilhamento de laboratórios e licenciamento via NIT) com os degraus da Escada do Capital (Centelha/FAPESB/FINEP $\to$ anjo $\to$ FNE Inovação) para reduzir a diluição societária dos fundadores antes de captar \textit{Venture Capital}.
  \espacoResolucao[2.8cm]
\end{enumerate}

\clearpage
\section*{Gabarito Comentado, Espelho Dissertativo e Diagnóstico: Simulado II}

\begin{boxGabarito}[Gabarito Comentado das Questões Objetivas (1 a 8)]
  \begin{enumerate}[leftmargin=1.5em, itemsep=4pt]
    \item \textbf{Gabarito: (B).} Limitar o desembolso a R\$\,15.000,00 próprios corresponde à perda aceitável (\textit{affordable loss}); firmar pré-compromissos com cooperativas para co-testar o produto corresponde à colcha de retalhos (\textit{crazy quilt}).
    \item \textbf{Gabarito: (B).} $FCA_0 = -90.000$; $FCA_1 = -200.000$; $FCA_2 = -240.000$ (fundo do Vale da Morte); $FCA_3 = -120.000$; $FCA_4 = +40.000$. $\text{Payback} = 3 + 120.000 / 160.000 = 3{,}75\text{ anos}$.
    \item \textbf{Gabarito: (A).} $\text{TAM} = 50.000 \times 2.400 = \text{R\$\,120.000.000,00}$ (R\$\,120,0 milhões). $\text{SAM} = 120.000.000 \times 0{,}30 = \text{R\$\,36.000.000,00}$ (R\$\,36,0 milhões). $\text{SOM}_{\text{BU}} = 450 \times 2.400 = \text{R\$\,1.080.000,00}$ (R\$\,1,08 milhão).
    \item \textbf{Gabarito: (B).} Processos industriais inéditos são protegidos por Patente de Invenção; o código proprietário agregado pode ser protegido por Registro de Software ou Segredo Industrial, observando os termos da licença aberta utilizada na base.
    \item \textbf{Gabarito: (B).} $\text{CAC} = 15.000 / 10 = \text{R\$\,1.500,00}$. $\text{MC}_{\text{mensal}} = 400 \times 0{,}75 = \text{R\$\,300,00}$. Retenção $= 1 / 0{,}05 = 20\text{ meses}$. $\text{LTV} = 300 \times 20 = \text{R\$\,6.000,00}$. $\text{LTV/CAC} = 6.000 / 1.500 = 4{,}00$.
    \item \textbf{Gabarito: (A).} $\text{MRR} = 50 \times 400 = \text{R\$\,20.000,00}$. $\text{Net Burn} = 60.000 - 20.000 = \text{R\$\,40.000,00/mês}$. $\text{Runway} = 320.000 / 40.000 = 8{,}0\text{ meses}$.
    \item \textbf{Gabarito: (B).} Subvenções não exigem participação societária nem devolução do principal, atuando exatamente na travessia TRL 3–6; crédito FNE é reembolsável (não é \textit{equity}).
    \item \textbf{Gabarito: (B).} $R_{\text{roy}} = 1.000.000 \times 0{,}05 = \text{R\$\,50.000,00}$. Parcela dos inventores $= 50.000 \times 0{,}30 = \text{R\$\,15.000,00}$. Parcela da ICT/NIT $= 50.000 - 15.000 = \text{R\$\,35.000,00}$.
  \end{enumerate}
\end{boxGabarito}

\begin{boxGabarito}[Espelho Dissertativo: Questões 9 e 10]
  \begin{itemize}[leftmargin=1.4em, itemsep=4pt]
    \item \textbf{Questão 9 (Critérios de Correção):}
    (a) \textit{Effectuation}: mobilizar \textit{Pássaro na mão} (uso do laboratório universitário para protótipos) e \textit{Colcha de retalhos} (cooperativa cacaueira cede amêndoas e galpão para validar TRL 5–6 com perda aceitável definida).
    (b) \textit{SOM bottom-up}: estimar quantas unidades a oficina local consegue montar e instalar por ano multiplicado pelo preço/assinatura anual, confrontando com o SAM do Sul da Bahia.
    (c) \textit{Fontes de Receita / BMC}: substituir venda de ativo por assinatura/serviço (HaaS ou cobrança por lote secado), eliminando a barreira de R\$\,35.000,00 de CAPEX inicial do produtor.
    \item \textbf{Questão 10 (Critérios de Correção):}
    Deve articular que o compartilhamento de infraestrutura da ICT (Lei nº 13.243/2016) e o licenciamento com *royalties* variáveis via NIT reduzem o \textit{Burn Rate} inicial. Ao combinar isso com subvenções não dilutivas (Centelha, FAPESB, FINEP) para elevar a maturidade de TRL 3 a TRL 6–7 e gerar MRR inicial, a \textit{spin-off} aumenta seu \textit{valuation pre-money} antes de buscar investidor-anjo ou \textit{Venture Capital} (e pode usar crédito FNE Inovação para planta produtiva), preservando maior participação dos fundadores.
  \end{itemize}
\end{boxGabarito}

\section*{Tabela de Diagnóstico e Plano de Revisão (Unidade II)}

\begin{center}
\renewcommand{\arraystretch}{1.24}
\small
\begin{tabularx}{\textwidth}{@{}c l >{\raggedright\arraybackslash}X l@{}}
  \toprule
  \rowcolor{etIndigo}
  \textcolor{white}{\bfseries Questões} & \textcolor{white}{\bfseries Capítulo} & \textcolor{white}{\bfseries Competência Avaliada} & \textcolor{white}{\bfseries Acerto (\checkmark / $\times$)} \\
  \midrule
  1 e 2 & Cap. 7 & \textit{Effectuation}, Curva J, Vale da Morte e \textit{payback} & \rule{1.8cm}{0.4pt} \\
  \rowcolor{bgGray}
  3 e 4 & Cap. 8 & TAM/SAM/SOM (\textit{top-down}/\textit{bottom-up}) e PI/FTO & \rule{1.8cm}{0.4pt} \\
  5 & Cap. 9 & BMC, \textit{Lean Startup}, CAC, LTV e razão LTV/CAC & \rule{1.8cm}{0.4pt} \\
  \rowcolor{bgGray}
  6 e 7 & Cap. 10 & Escada do Capital, fomento regional, MRR, \textit{Burn} e \textit{Runway} & \rule{1.8cm}{0.4pt} \\
  8 & Cap. 11 & \textit{Open Innovation}, \textit{spin-off}, NIT e \textit{royalties} (Lei 13.243/2016) & \rule{1.8cm}{0.4pt} \\
  \rowcolor{bgTeal}
  9 e 10 & Síntese & Integração estratégica e territorial para a AV3 & \rule{1.8cm}{0.4pt} \\
  \bottomrule
\end{tabularx}
\end{center}
